\section*{Введение}
\addcontentsline{toc}{section}{Введение}

В данной курсовой работе рассматривается система управления модулем поворота промышленного робота с аналоговым ПИ-регулятором. Основной задачей является подбор параметров обратной связи, обеспечивающих устойчивую работу системы и удовлетворение заданным ограничениям на управляющее воздействие и механические нагрузки.

Для исследуемой системы требуется обеспечить поворот механизма на заданный угол $\varphi^* = \pi$, при этом величина управляющего сигнала не должна превышать допустимого значения $u < 220$ В, а момент на выходном валу редуктора должен удовлетворять условию
\[
|M_{\text{м}}|_{\max} < 150 \text{ Нм}.
\]

Для определения допустимых параметров регулятора необходимо построить область устойчивости системы управления в пространстве коэффициентов обратной связи. Исследование устойчивости проводится с использованием критериев Стодолы и Рауса--Гурвица. Дополнительно устойчивость системы в отдельных точках проверяется с помощью критерия Михайлова.

После построения области устойчивости определяется область параметров, удовлетворяющая ограничениям на крутящий момент и управляющий сигнал. В найденной области строятся линии равной длительности переходных процессов и определяются параметры регулятора, обеспечивающие минимальное время регулирования.

Также в работе рассматривается задача оптимального управления. Для этого используется интегральный функционал качества, учитывающий как ошибку регулирования, так и величину управляющего воздействия. На основе полученных результатов исследуются переходные процессы при различных значениях коэффициентов обратной связи.

В заключительной части работы строятся графики перемещения выходного звена $\varphi_{\text{м}}$, сигнала рассогласования $\psi_{\text{р}}$, упругой деформации редуктора $\varphi_{\text{р}}/i - \varphi_{\text{м}}$, крутящего момента $M_{\text{м}}$, приведенного момента двигателя $M_{\text{дпр}}$ и сигнала управления $u$ как функций времени.
